% =============================================================
% Section 4. Finite Element Exterior Calculus: W^0 Lift and P Projection
%
% Character: comparative review chapter. We introduce no new theory.
% =============================================================

% --- Opening hook (continuing from Section 3) ---
The technical core of Robinson's open problem---a sheaf-categorically
intrinsic bandwidth---requires lifting Hilbert-complex structure,
Hodge decomposition, and bounded cochain projection from the
continuous setting of finite element exterior calculus (FEEC) to
cellular sheaves of Hilbert spaces. We review the necessary FEEC
machinery in this section. The exposition is comparative: four
schools---mimetic discretization \cite{BochevHyman2006}, finite
element systems \cite{Christiansen2010}, FEEC Hilbert complexes
\cite{AFW2010}, and dagger Galerkin \cite{LahtinenStenvall2020}---
have developed the same family of constructions in disjoint
vocabularies over three decades. \emph{This section is a
comparative review; we do not introduce new theory.} Its purpose is
disambiguation and a side-by-side mapping that makes the
sheaf-side interface of \cref{sec:sheafbw} statable.

% =============================================================
\subsection{Bochev--Hyman reduction and reconstruction}\label{subsec:feec-bh}
% =============================================================

The mimetic-discretization framework of Bochev--Hyman
\cite{BochevHyman2006} organizes the connection between continuous
differential forms and discrete cochains around a pair of maps and a
commuting diagram.

\begin{definition}[Reduction and Whitney reconstruction]\label{def:bh-RI}
Let $\mathcal{T}$ be a cellular complex and $\Lambda^k(\Omega)$ the
space of differential $k$-forms on a domain $\Omega \supset
\lvert \mathcal{T} \rvert$. The \emph{reduction} (or de Rham map)
\[
  R : \Lambda^k(\Omega) \longrightarrow C^k(\mathcal{T}),
  \qquad
  R(\omega)(\sigma) := \int_\sigma \omega
  \quad\text{for $k$-cells } \sigma \in \mathcal{T},
\]
is the integration cochain map \cite[\S 1]{BochevHyman2006}, and the
\emph{Whitney reconstruction}
\[
  I_W : C^k(\mathcal{T}) \longrightarrow V_W^k \subset \Lambda^k(\Omega),
\]
sends a cochain to the unique Whitney $k$-form interpolating it on
each cell of $\mathcal{T}$
\cite{Whitney1957,Bossavit1988}.\footnote{Bochev--Hyman~%
\cite{BochevHyman2006} originally write $R/I$ for the
reduction/reconstruction pair. To avoid clash with the integration
map $I$ used in \cref{sec:three_instances} for the
sampling/integration duality, we write $R/I_W$ throughout this paper,
with subscript $W$ indicating Whitney reconstruction.}
\end{definition}

\begin{proposition}[Commuting Stokes; cf.\ {\cite[\S 2]{BochevHyman2006}}]
\label{prop:commuting-stokes}
With $R$ and $I_W$ as above, $R \circ d = \delta \circ R$, where $d$
is the exterior derivative on $\Lambda^k(\Omega)$ and $\delta$ is the
cochain coboundary on $C^k(\mathcal{T})$. Equivalently, $R$ is a
chain map between the de Rham complex of $\Omega$ and the cochain
complex of $\mathcal{T}$. The natural inner product on
$C^k(\mathcal{T})$ induced by Whitney reconstruction is
$\langle c, c' \rangle_{C^k} := \langle I_W c, I_W c' \rangle_{L^2(\Omega)}$.
\end{proposition}

The pair $(R, I_W)$ behaves asymmetrically with respect to composition:

\begin{remark}[$I_W$ is not a projection]\label{rem:IW-not-projection}
On $V_W^k \subset \Lambda^k(\Omega)$ one has $R \circ I_W = \mathrm{id}_{C^k}$
\cite{Whitney1957}, but on the larger ambient $\Lambda^k(\Omega)$ the
composition $I_W \circ R$ is not the identity, and $I_W$ is not a
bounded $L^2$-projection onto $V_W^k$. The same observation is
recorded explicitly by Falk--Winther \cite{FalkWinther2014}: ``the
operator is not a projection since it is not equal to the identity on
the Whitney forms.'' The repair is the construction of a
\emph{bounded cochain projection} $\pi_h$ in
\cref{subsec:feec-afw}.
\end{remark}

% =============================================================
\subsection{Christiansen finite element systems}\label{subsec:feec-fes}
% =============================================================

The mimetic framework of \cref{subsec:feec-bh} acquires a contravariant
functor form in Christiansen's finite element system (FES) framework
\cite{Christiansen2010}.

\begin{definition}[Finite element system; {\cite[\S 2, Def.~2]{Christiansen2010}}]
\label{def:christiansen-fes}
Let $\mathcal{T}$ be a cellular complex. A \emph{finite element system}
on $\mathcal{T}$ assigns, to each $k \in \mathbb{N}$ and each cell
$T \in \mathcal{T}$, a vector space
$\mathcal{E}^k(T) \subseteq \widehat{X}^k(T)$
of \emph{differential $k$-elements}\footnote{We follow
\cite[\S 2, Def.~2]{Christiansen2010} with one notational change: we
write $\mathcal{E}^k(T)$ for what \cite{Christiansen2010} denotes
$A^k(T)$, since $A$ is reserved throughout this paper for the
sampling operator $A : V_B \to \C^M$ of \cref{eq:A-def}.}
such that
\begin{enumerate}
  \item $d : \mathcal{E}^k(T) \to \mathcal{E}^{k+1}(T)$ is well
        defined; and
  \item for every cell inclusion $i : T' \subseteq T$, the pullback
        $i^\star : \mathcal{E}^k(T) \to \mathcal{E}^k(T')$ is a
        morphism of vector spaces.
\end{enumerate}
\end{definition}

\Cref{def:christiansen-fes} exhibits an FES as a contravariant functor
from the poset category of $\mathcal{T}$ (objects: cells; morphisms:
inclusions) to $\Vect$, with $T \mapsto \mathcal{E}^k(T)$ and
inclusions sent to pullbacks. Christiansen and Rapetti
\cite{ChristiansenRapetti2025} state this categorical reading
explicitly: ``a FES is just a contravariant functor from the
(poset) category determined by the mesh, to the category of vector
spaces.'' Whitney $k$-forms are an instance:
\cite[Example~2.3]{Christiansen2010} identifies the case of the
trimmed polynomial complex of degree~$1$ with the Whitney form
construction of \cite{Whitney1957}, applied in computational
electromagnetics by \cite{Bossavit1988}.

The de Rham map of an FES recovers cohomology in the following sense.

\begin{proposition}[{\cite[Prop.~2.10]{Christiansen2010}}]\label{prop:fes-derham}
Under the compatibility hypotheses of
\cite[\S 2, Def.~3]{Christiansen2010}, the de Rham map from an FES to
the simplicial cochain complex of $\mathcal{T}$ induces an isomorphism
on cohomology.
\end{proposition}

% =============================================================
\subsection{AFW Hilbert complexes and bounded cochain projection}\label{subsec:feec-afw}
% =============================================================

The framework of \cref{subsec:feec-bh,subsec:feec-fes} acquires its
analytic completion in the FEEC Hilbert-complex framework of
Arnold--Falk--Winther \cite{AFW2006,AFW2010}, which extends the
algebraic structure to closed densely-defined operators between
Hilbert spaces.

\begin{definition}[Hilbert complex; {\cite[\S 3.1.3]{AFW2010},
\cite{BruningLesch1992}}]\label{def:afw-hilbert-complex}
A \emph{Hilbert complex} is a sequence $(W^k, d^k)_{k \ge 0}$ in
which each $W^k$ is a Hilbert space and each
$d^k : W^k \to W^{k+1}$ is a closed, densely defined linear operator
satisfying $\mathrm{range}(d^k) \subseteq \mathrm{dom}(d^{k+1})$ and
$d^{k+1} \circ d^k = 0$. The complex is \emph{closed} when, for each
$k$, the range of $d^k$ is closed in $W^{k+1}$.
\end{definition}

Closedness yields the standard Hodge decomposition. Writing
$(d^k)^*$ for the Hilbert adjoint of $d^k$, set
$\Delta_k := d^{k-1} (d^{k-1})^* + (d^k)^* d^k$.

\begin{theorem}[Hodge decomposition;
{\cite[\S 3.2]{AFW2010},\cite{BruningLesch1992}}]\label{thm:afw-hodge}
For a closed Hilbert complex $(W^k, d^k)$,
\[
  W^k \;=\; \overline{\mathrm{range}(d^{k-1})}
  \;\oplus\; \mathrm{ker}(\Delta_k)
  \;\oplus\; \overline{\mathrm{range}((d^k)^*)},
\]
the three summands being mutually orthogonal. The harmonic space
$\mathcal{H}^k := \mathrm{ker}(\Delta_k)$ is canonically isomorphic
to the cohomology $H^k$ of the complex.
\end{theorem}

The discrete-side structure that connects $W^k$ to a finite-dimensional
subcomplex $W_h^k$ is the \emph{bounded cochain projection}.

\begin{definition}[Bounded cochain projection;
{\cite[\S 3.1.1]{AFW2010},\cite{FalkWinther2014}}]\label{def:bounded-pi}
A family of operators $\pi_h^k : W^k \to W_h^k$ is a
\emph{bounded cochain projection} when, for each $k$, it satisfies
\begin{enumerate}
  \item $\pi_h^k \circ \iota_h = \mathrm{id}_{W_h^k}$, where
        $\iota_h : W_h^k \hookrightarrow W^k$ is the canonical
        inclusion (true projection);
  \item $d^k \circ \pi_h^k = \pi_h^{k+1} \circ d^k$ on
        $\mathrm{dom}(d^k)$ (commutation with the coboundary);
  \item $\lVert \pi_h^k \rVert$ is uniformly bounded with respect to
        the mesh parameter $h$.
\end{enumerate}
\end{definition}

\begin{theorem}[{\cite[Thm.~3.6]{AFW2010}}]\label{thm:afw-cohomology}
If $(W_h^k, d^k)$ is a subcomplex of a closed Hilbert complex
$(W^k, d^k)$ admitting a bounded cochain projection, then the
cohomology of the subcomplex is canonically isomorphic to that of
the ambient complex, and the harmonic spaces have equal dimension.
\end{theorem}

Falk--Winther \cite{FalkWinther2014} construct a \emph{local}
bounded cochain projection meeting all three conditions of
\cref{def:bounded-pi} on simplicial subcomplexes; the construction
repairs the obstruction in \cref{rem:IW-not-projection} by replacing
$I_W$ with a smoothed extension that retains the commutation with
$d$ while becoming a true projection. Christiansen
\cite[\S 4]{Christiansen2010} gives a related family of
quasi-interpolators.

% =============================================================
\subsection{Lahtinen--Stenvall dagger mono form of Galerkin}\label{subsec:feec-dagger}
% =============================================================

A complementary categorical reading is provided by Lahtinen and
Stenvall \cite{LahtinenStenvall2020}, who characterize Galerkin
discretization as a morphism in the dagger category
$\Hilb$ of Hilbert spaces with bounded linear maps.

A \emph{dagger category} \cite{HeunenVicary2019} is a category
equipped with an involutive contravariant identity-on-objects functor
$\dagger$; in $\Hilb$, $f^\dagger$ is the Hilbert adjoint of
$f : H_1 \to H_2$. A \emph{dagger mono} is a morphism
$\iota : H_h \to H$ with $\iota^\dagger \circ \iota = \mathrm{id}_{H_h}$,
equivalently an isometric embedding.

\begin{proposition}[Galerkin as dagger mono;
{\cite{LahtinenStenvall2020}}]\label{prop:galerkin-dagger}
Discretization of a finite element method, viewed as a morphism in
$\Hilb$, is a dagger mono $\iota_h : V_h \hookrightarrow V$ from a
finite-dimensional discrete domain into the continuous Hilbert
space. The Galerkin projection $V \to V_h$ is the dagger
$\iota_h^\dagger$ of this embedding.
\end{proposition}

\Cref{prop:galerkin-dagger} aligns the FEEC framework with the
dagger-categorical formulation: $\iota_h$ is a dagger mono in
$\Hilb$ in the same sense as the inclusion of the Whitney form
subspace $V_W^k \hookrightarrow \Lambda^k(\Omega)$, and the Galerkin
projection is its Hilbert adjoint. The bounded cochain projection
$\pi_h$ of \cref{def:bounded-pi} is not, in general, a dagger
projection (i.e., it need not equal $\iota_h^\dagger$), since the
local construction of \cite{FalkWinther2014} sacrifices
self-adjointness in exchange for locality and uniform boundedness.

% =============================================================
\subsection{Four-school comparison table}\label{subsec:feec-comparison}
% =============================================================

The four schools of \cref{subsec:feec-bh,subsec:feec-fes,subsec:feec-afw,subsec:feec-dagger}
develop the same family of constructions in disjoint vocabularies.
\Cref{tab:feec-four-schools} records the side-by-side mapping.

\begin{table}[h]
  \centering
  \small
  \begin{tabular}{@{}llll@{}}
    \toprule
    Construction & Bochev--Hyman \cite{BochevHyman2006} &
       Christiansen \cite{Christiansen2010} &
       AFW \cite{AFW2010} \\
    \midrule
    Continuous space &
       $\Lambda^k(\Omega)$ & $\widehat{X}^k(T)$ & $W^k$ \\
    Discrete space &
       $C^k(\mathcal{T})$ & $\mathcal{E}^k(T)$ & $W_h^k$ \\
    Coboundary & $\delta$ on $C^k$ & $d$ on $\mathcal{E}^k$ &
       $d^k$ on $W^k$ \\
    Lift to ambient &
       $I_W$ (Whitney) & FES inclusions \cite[Ex.~2.3]{Christiansen2010} &
       $\iota_h : W_h^k \hookrightarrow W^k$ \\
    Reduction to discrete &
       $R$ (de Rham map) & pullbacks $i^\star$ & $\pi_h^k$
       (bounded cochain projection) \\
    Commutation &
       $R \circ d = \delta \circ R$ & functoriality of FES &
       $d^k \circ \pi_h^k = \pi_h^{k+1} \circ d^k$ \\
    Cohomology agreement &
       (implicit via Whitney 1957) & \cite[Prop.~2.10]{Christiansen2010} &
       \cite[Thm.~3.6]{AFW2010} \\
    \bottomrule
  \end{tabular}
  \caption{The mimetic, FES, and AFW frameworks for the
    same family of constructions, mapped column by column. The
    Lahtinen--Stenvall framework \cite{LahtinenStenvall2020} reads
    the inclusion row as a dagger mono in
    $\Hilb$.}\label{tab:feec-four-schools}
\end{table}

\begin{remark}[Common scope]\label{rem:feec-scope}
Each of the four schools establishes its constructions for a single
Hilbert complex (continuous side) related to a single subcomplex
(discrete side). None of the four schools attaches the constructions
to a sheaf over a cell complex with non-trivial restriction maps
between stalks. The transition to the sheaf side is taken in
\cref{sec:sheafbw}.
\end{remark}

% =============================================================
\subsection{Robinson's open problem: technical interface}\label{subsec:feec-robinson-interface}
% =============================================================

Robinson \cite[\S 6]{Robinson2014} writes that ``If $F$ and $S$ are
sheaves of Hilbert spaces then their cochain complexes are Hilbert
complexes, and they have a Hodge decomposition\ldots\ There may be
a way to discuss the concept of \emph{bandwidth} for general sheaves
of Hilbert spaces.'' \Cref{thm:triangle-equiv,def:sheafbw-candidate}
narrowed this to the two-clause problem
\textup{(R1)}--\textup{(R2)} of \cref{subsec:tri-narrowing}.
\Cref{sec:sheafbw} answers it in three steps, each of which uses one
ingredient of the FEEC machinery reviewed above.

\begin{enumerate}
  \item[\textbf{Step 1:}] \emph{Hilbert complex structure on cellular
        sheaves of Hilbert spaces.} Each stalk
        $\mathscr{A}(\sigma)$ becomes a Hilbert space, restriction
        maps become bounded operators, and the cochain spaces
        $C^k(\mathscr{A})$ assemble into a closed Hilbert complex
        in the sense of \cref{def:afw-hilbert-complex}. In the
        finite-dimensional setting of \cref{sec:sheafbw}, closedness
        of cochain ranges is automatic
        (\cref{thm:closed-hilbert-complex}). This step lifts
        \cref{def:afw-hilbert-complex} from a single chain to a
        sheaf-valued chain.
  \item[\textbf{Step 2:}] \emph{Hodge decomposition stalk-wise with
        restriction commutation.} \Cref{thm:afw-hodge} applied to
        the sheaf-cochain Hilbert complex of Step~1 yields a
        decomposition of $C^k(\mathscr{A})$ into exact, harmonic,
        and co-exact summands; the restriction maps of $\mathscr{A}$
        commute with the decomposition by functoriality. The
        sheaf Laplacian of \cite{HansenGhrist2019} is the
        finite-dimensional specialization.
  \item[\textbf{Step 3:}] \emph{Categorically intrinsic sheaf
        bandwidth.} \Cref{def:bounded-pi}'s view of
        well-conditioning as the boundedness of $\pi_h$ is replaced
        on the sheaf side by a spectral-gap ratio of the sheaf
        coboundary $\delta^0$, namely
        $\sheafBW := \sqrt{\gamma_+(\delta^0) /
        \Gamma_+(\delta^0)}$, defined intrinsically on the cellular
        sheaf data of $\mathscr{A}$. This step answers
        clauses~\textup{(R1)} and \textup{(R2)} of
        \cref{subsec:tri-narrowing}; it is constructed in
        \cref{sec:sheafbw}.
\end{enumerate}

Steps~1 and~2 are direct transcriptions of \cite{AFW2010}
\cref{def:afw-hilbert-complex,thm:afw-hodge} into the sheaf setting
and contribute no new theory. Step~3 is the categorically intrinsic
bandwidth definition; it is the locus of the new content of this
paper, completed in \cref{sec:sheafbw}.

% --- Closing hook (to Section 5) ---
\Cref{def:afw-hilbert-complex,thm:afw-hodge,def:bounded-pi}, together
with the FES contravariant functor form
\cref{def:christiansen-fes,prop:fes-derham} and the dagger
formulation \cref{prop:galerkin-dagger}, supply the algebraic
machinery: closed Hilbert complexes, Hodge decompositions, and
bounded cochain projections, with the categorical structures that
make their sheaf-side transcription unambiguous. What is missing on
the sheaf side---a categorically intrinsic bandwidth answering
clauses~\textup{(R1)} and~\textup{(R2)} of
\cref{subsec:tri-narrowing}---is constructed in
\cref{sec:sheafbw}.
